\chapter{Einleitung}
In der modernen Softwareentwicklung vollzieht sich ein fundamentaler Wandel, der weit über technologische Neuerungen hinausgeht.
\DO, oft als bloße Methodik oder Werkzeugsatz missverstanden, ist in seinem Kern ein kultureller und organisatorischer Transformationsprozess.
Es geht nicht mehr nur darum, \emph{was} entwickelt wird, sondern vor allem darum, \emph{wie} Organisationen zusammenarbeiten, um Innovationen schnell, zuverlässig und wertschöpfend an den Markt zu bringen.

Die Motivation für die Auseinandersetzung mit \DO entspringt dem steigenden Druck auf Unternehmen, in einem hochvolatilen Marktumfeld zu bestehen.
Die Fähigkeit, neue Produktfeatures in kurzen Zyklen zu liefern und die Softwarequalität kontinuierlich zu verbessern, ist zu einem entscheidenden Wettbewerbsvorteil geworden~\cite[1]{mak25}.

\DO verspricht genau das: die Beschleunigung von Lieferzyklen und die Steigerung der Stabilität durch die Überwindung traditioneller Differenzen zwischen Entwicklungs"= (Dev) und Betriebsteams (Ops)~\cite[284]{kar24}.
Die zentrale Problemstellung liegt jedoch darin, dass dieser Wandel nicht durch die alleinige Einführung von Automatisierungswerkzeugen erreicht wird.
Vielmehr handelt es sich um einen tiefgreifenden soziokulturellen Prozess, dessen Hauptziel der Aufbau einer kollaborativen Kultur ist~\cite[4]{luz18}.
Die Vernachlässigung dieser sozialen Dimension ist gleichzeitig eine der häufigsten Ursachen für das Scheitern von \DO-Initiativen.

Aus dieser Problemstellung leitet sich die zentrale Forschungsfrage dieser Arbeit ab: Welche sozialen und kulturellen Faktoren fördern oder hemmen die \DO-Transformation in Software-Unternehmen?
Das Ziel ist es, diese soziokulturellen Erfolgs"= und Hemmnisfaktoren auf Basis der existierenden wissenschaftlichen Literatur zu identifizieren und zu analysieren.
Es soll ein klares Bild davon gezeichnet werden, warum die menschliche und organisatorische Dimension den Dreh"= und Angelpunkt einer erfolgreichen \DO-Transformation darstellt.

Die Arbeit ist wie folgt strukturiert: \autoref{ch:theoretical-framework} legt den theoretischen Rahmen dar und definiert \DO als soziotechnisches Phänomen.
\autoref{ch:social-cultural-factors} analysiert die zentralen sozialen und kulturellen Erfolgsfaktoren, die eine \DO-Transformation begünstigen.
\autoref{ch:barriers} beleuchtet hingegen die Hemmnisse und Widerstände, die in traditionellen Organisationsstrukturen auftreten.
In \autoref{ch:discussion} werden die gewonnenen Erkenntnisse diskutiert und in den breiteren Kontext des organisationalen Wandels und der modernen Arbeitswelt eingeordnet.
Das abschließende \autoref{ch:conclusion} fasst die Hauptergebnisse zusammen, beantwortet die Forschungsfrage und gibt einen Ausblick auf zukünftige Forschungsfelder.